\section{Implementation and Model Alignment}
\label{sec:implementation}
\subsection{Mooncake prototype}
The prototype adds a native C++ AegisController to Mooncake Store,
Application/Router services, and a vLLM connector gate. Request-specific
cache salts remain stable across retries but prevent unrelated requests
from sharing physical KV. The gate seals completed forward buffers,
commits successful tier copies, and drains pending work before cleanup.
Recovery loads local HBM into vLLM slots before fetching remaining Store
keys. Section~\ref{sec:physical} reports physical validation.

\subsection{MoonKVSim projection}
MoonKVSim uses one global nanosecond clock, per-worker schedulers,
HBM/DRAM ledgers and shared fluid-rate resources. KV becomes readable at
forward-step completion. Transfers reserve destinations and track source
lifetime, service, completion and commitment. Failed GPU memory becomes
unavailable and returns empty after repair.

The \texttt{mooncake\_prompt\_v1} projection models all-chunk candidates,
4,096-token prompt sealing, request salt, rendezvous routing, the shared
ordinal barrier, independent tier decisions, recovery copying and cleanup.
Its AIMD interval is 1 s versus 5 s in the physical pilot. Recovery reserves
separate primary memory and waits on the bound target when capacity is
insufficient. Prior-model executions are excluded from every revised result.

\subsection{What is and is not reproduced}
Table~\ref{tab:alignment} defines the evidence boundary. The frozen workload
contains exact prompt token IDs and prescribed output counts, not generated
continuations or KV tensors. We therefore model prompt protection and replay
unprotected output positions. The real connector's 256-token decode protection
is described as implementation behavior and receives no simulated recovery
credit. The simulator's 64-token physical allocation unit can overallocate a
partial tail relative to the prototype's 16-token vLLM pages; it does not pad
the workload's token sequence.

\begin{table*}[t]
\centering
\caption{Alignment with the Mooncake reference and remaining model boundaries.
Conformance of metadata does not calibrate device timing.}
\label{tab:alignment}
\begin{tabular}{@{}p{.21\textwidth}p{.37\textwidth}p{.37\textwidth}@{}}
\toprule
Aspect & Mooncake prototype & Evaluated projection \\
\midrule
Tier ownership & Source-owned independent HBM/DRAM & Same admission, commit and lifetime contract \\
Chunk visibility & Completed forward buffers; shared ordinal & Same atomic visibility and ordinal barrier \\
Recovery & Local copy locking; visibility lease & Explicit copy; source pinned until load completes \\
AIMD interval & 5 s in the accepted physical pilot & 1 s in the frozen simulation \\
Application signal & Causal GBDT plus calibrated residual CPM & Prescribed ideal durations; zero-Slack ablation \\
Protected state & Prompt and sealed decode KV tensors & Exact prompt token intervals; decode replay only \\
Allocation granularity & 16-token vLLM pages & 64-token source allocation units \\
\bottomrule
\end{tabular}
\end{table*}

The native HBM free path hides recovery keys and frees slots; the Controller
does not consult Worker leases when selecting pressure victims. The local
copy is locked, but eviction between activation and load may require fallback.
The simulator excludes pinned victims throughout materialization, a stronger
retention assumption. The differential checks do not validate this race.
AegisController metadata is in-memory; integration with Store HA durability
is outside the demonstrated prototype.

\subsection{Verification}
A harness compiles the exact native Controller and compares 260 policy
cases: tier-local gain, independence, pending invisibility, Slack expiry
and AIMD updates. Six further sequences exercise the reference ordinal
finalizer, including all-bypass retries. These checks establish metadata
conformance, not concurrent data-plane equivalence.

Overlapping regression suites pass 127 and 184 tests. Each execution audits
resources, policy, faults and cleanup; independent event reconstruction checks
workflow times. Hashes bind execution source, workload and analysis.
